\startSection
\selectlanguage{german}
\sectionTitle{Die Deutsche Nationalhymne}

\begin{question}{Von welchem Komponisten stammt die Melodie, die heute mit der deutschen Nationalhymne verbunden ist?}
  \choice{Ludwig van Beethoven}
  \choice{Johann Sebastian Bach}
  \choice[correct]{Joseph Haydn}
  \choice{Richard Wagner}
  \choice{Johannes Brahms}
\end{question}

\begin{question}{Wer schrieb den Text des \enquote{Deutschlandliedes}?}
  \choice{Johann Wolfgang von Goethe}
  \choice{Friedrich Schiller}
  \choice{Heinrich Heine}
  \choice[correct]{August Heinrich Hoffmann von Fallersleben}
  \choice{Theodor Körner}
\end{question}

\begin{question}{Welche Strophe des \enquote{Deutschlandliedes} wird heute als deutsche Nationalhymne gesungen?}
  \choice{Die erste Strophe}
  \choice{Die zweite Strophe}
  \choice[correct]{Die dritte Strophe}
  \choice{Alle drei Strophen}
  \choice{Die erste und dritte Strophe}
\end{question}

\begin{question}{Mit welchen Worten beginnt die heute verwendete Strophe?}
  \choice{\enquote{Deutschland, Deutschland über alles}}
  \choice{\enquote{Deutsche Frauen, deutsche Treue}}
  \choice[correct]{\enquote{Einigkeit und Recht und Freiheit}}
  \choice{\enquote{Blüh im Glanze dieses Glückes}}
  \choice{\enquote{Von der Maas bis an die Memel}}
\end{question}

\begin{question}{In welchem Jahr wurde der Text des \enquote{Deutschlandliedes} ursprünglich geschrieben?}
  \choice{1797}
  \choice{1815}
  \choice[correct]{1841}
  \choice{1871}
  \choice{1918}
\end{question}

\begin{question}{Welche politische Idee steht im Zentrum der heute gesungenen dritten Strophe?}
  \choice{Die Wiederherstellung einer absoluten Monarchie}
  \choice{Die Ausweitung kolonialer Herrschaft}
  \choice{Die Abschaffung aller Bundesländer}
  \choice[correct]{Die Verbindung von Einheit, Recht und Freiheit}
  \choice{Die Gründung einer europäischen Armee}
\end{question}
\shuffleSection
\renderSection